\subsection{更换基环。}

本节中，我们用 A、B、A′、B′ 表示四个环，用 h 与 h′ 分别表示 A 到 A′ 与 B 到 B′ 的同态，用 G 表示 (A, B)-双模，用 G′ 表示 (A′, B′)-双模，并用 u 表示 G 的底加法群到 G′ 的底加法群的一个同态，它满足

\[
u (a g b) = h (a) u (g) h ^ {\prime} (b) \quad (a \in \mathrm{A}, g \in \mathrm{G}, b \in \mathrm{B}).\tag{11}
\]

设 E（相应地 F）为左 A-模（相应地右 B-模）。回顾（第 III 章，第 2 版，附录 II，第 10 小节）：若把 A′（相应地 B′）看作右 A-模（相应地左 B-模），则张量积 \(E' = A' \otimes_{A} E\) （相应地 \(F' = F \otimes_{B} B'\) ）带有由下式定义的左 A′-模（相应地右 B′-模）结构

\[
\begin{array}{c c} a _ {1} ^ {\prime} (a ^ {\prime} \otimes x) = (a _ {1} ^ {\prime} a ^ {\prime}) \otimes x & (a ^ {\prime}, a _ {1} ^ {\prime} \in A ^ {\prime}, x \in E) \\ (\text { resp. } (y \otimes b ^ {\prime}) b _ {1} ^ {\prime} = y \otimes (b ^ {\prime} b _ {1} ^ {\prime}) & (b ^ {\prime}, b _ {1} ^ {\prime} \in B ^ {\prime}, y \in F)). \end{array}\tag{12}
\]

命题 1. --- 设 E 为左 A-模，F 为右 B-模；令 \(\mathbf{E}^{\prime} = \mathbf{A}^{\prime}\otimes_{\mathbf{A}}\mathbf{E}\) 与 \(\mathbf{F}^{\prime} = \mathbf{F}\otimes_{\mathbf{B}}\mathbf{B}^{\prime}\) 。对每一个双线性映射 \(\Phi\) （从 \(\mathbf{E}\times \mathbf{F}\) 到 G 中），存在且仅存在一个双线性映射 \(\Phi^{\prime}\) （从 \(\mathbf{E}^{\prime}\times \mathbf{F}^{\prime}\) 到 G′ 中），使得有

\[
\Phi^ {\prime} (a ^ {\prime} \otimes x, y \otimes b ^ {\prime}) = a ^ {\prime}. u (\Phi (x, y)). b ^ {\prime}\tag{13}
\]

对一切 \(a' \in \mathrm{A}'\) 、\(b' \in \mathrm{B}'\) 、\(x \in \mathrm{E}\) 、\(y \in \mathrm{F}\) 。

\(\Phi'\) 的唯一性由元素 \(a' \otimes x\) 与 \(y \otimes b'\) 分别生成 E′ 与 F′ 这一事实得出。为证其存在性，考虑映射

\[
m: (a ^ {\prime}, x, y, b ^ {\prime}) \rightarrow a ^ {\prime}. u (\Phi (x, y)). b ^ {\prime}
\]

它把 A′ × E × F × B′ 映到 G 中；它显然是 Z-重线性的，且满足

\[
\begin{array}{r} m (a ^ {\prime}, a x, y, b ^ {\prime}) = m (a ^ {\prime} h (a), x, y, b ^ {\prime}) \\ m (a ^ {\prime}, x, y b, b ^ {\prime}) = m (a ^ {\prime}, x, y, h ^ {\prime} (b) b ^ {\prime}) \\ (a \in \mathrm{A}, b \in \mathrm{B}, a ^ {\prime} \in \mathrm{A} ^ {\prime}, b ^ {\prime} \in \mathrm{B} ^ {\prime}, x \in \mathrm{E}, y \in \mathrm{F}). \end{array}
\]

因此存在 E′ × F′ 到 G′ 中的、满足 (13) 的一个 Z-双线性映射 \(\Phi'\) （第 III 章，第 2 版，附录 II，第 1 小节，命题 2）。这一关系以及由 (12) 给出的 E′ 与 F′ 的模结构定义表明 \(\Phi'\) 是双线性的，证毕。

在命题 1 的假设与记号下，我们现在来研究与 \(\Phi\) 及 \(\Phi'\) 相伴的线性映射（第 1 小节，定义 2）。为此，我们先定义 \(\mathfrak{L}_{A}(\mathrm{E}, \mathrm{G})\) 到 \(\mathfrak{L}_{A'}(\mathrm{E}', \mathrm{G}')\) 中的一个典范同态。对每个 \(\nu \in \mathfrak{L}_{\Lambda}(\mathrm{E}, \mathrm{G})\) ，映射 \((a', x) \to a'. u(\nu(x))\) （从 \(\mathrm{A}' \times \mathrm{E}\) 到 \(\mathrm{G}'\) 中）是 Z-双线性的，且鉴于 (11)，它把 \((a'h(a), x)\) 与 \((a', ax)\) （ \(a \in \mathrm{A}\) ）映到 \(\mathrm{G}'\) 的同一元素上；因而它定义了（第 III 章，第 2 版，

附录 II，第 1 与 10 小节）一个映射 \(k(\nu)\) （从 \(\mathbf{E}' = \mathbf{A}' \otimes_{\mathbf{A}} \mathbf{E}\) 到 \(\mathbf{G}'\) 中），使得 \(k(\nu)(a' \otimes x) = a'. u(\nu(x))\) ，且由 (12)，它是 A′-线性的。此外，由 (12) 立即推出，映射 \(\nu \to k(\nu)\) （从 \(\mathfrak{L}_{\mathbf{A}}(\mathbf{E}, \mathbf{G})\) 到 \(\mathfrak{L}_{\mathbf{A}'}(\mathbf{E}', \mathbf{G}')\) 中）满足 \(k(\nu b) = k(\nu) h'(b)\) （对一切 \(b \in B\)）。用 \(i\) 表示典范映射 \(y \to y \otimes 1\) （从 F 到 \(\mathbf{F}'\) 中）。则图表

\[
\begin{array}{c c c} \mathrm{F} & \xrightarrow {d _ {\Phi}} & \mathcal {L} _ {A} (\mathrm{E}, \mathrm{G}) \\ \downarrow i & & \downarrow k \\ \mathrm{F} ^ {\prime} & \xrightarrow {d _ {\Phi^ {\prime}}} & \mathcal {L} _ {A^{\prime}} (\mathrm{E} ^ {\prime}, \mathrm{G} ^ {\prime}) \end{array}\tag{14}
\]

（其中 \(d_{\Phi}\) 与 \(d_{\Phi'}\) 表示 \(\Phi\) 与 \(\Phi'\) 的右相伴线性映射）是交换的。事实上，对 \(x \in \mathrm{E}, y \in \mathrm{F}\) 与 \(a' \in \mathrm{A}'\) ，有 \(d_{\Phi'}(i(y))(a' \otimes x) = \Phi'(a' \otimes x, y \otimes 1) = a'. u(\Phi(x, y)) = a'. u(d_{\Phi}(y)(x))\) ，也就是说 \(d_{\Phi'}(i(y))(a' \otimes x) = k(d_{\Phi}(y))(a' \otimes x)\) 。线性映射 \(s_{\Phi}\) 与 \(s_{\Phi'}\) （它们左相伴于 \(\Phi\) 与 \(\Phi'\) ）有类似的交换关系。

命题 2. --- 设 B 与 B′ 分别配备了反自同构 J 与 I，使得

\[
h ^ {\prime} (b ^ {J}) = h ^ {\prime} (b) ^ {I} \quad \text { for all } b \in \mathrm{B}.\tag{15}
\]

设 E 为左 A-模，F 为左 B-模；令 \(E' = A' \otimes_{A} E\) 与 \(F' = B' \otimes_{B} F\) 。对每一个（关于 J 的）半双线性映射 \(\Phi\) （从 \(E \times F\) 到 G 中），存在且仅存在一个（关于 I 的）半双线性映射 \(\Phi'\) （从 \(E' \times F'\) 到 \(G'\) 中），使得

\[
\Phi^ {\prime} (a ^ {\prime} \otimes x, b ^ {\prime} \otimes y) = a ^ {\prime}. u (\Phi (x, y)). b ^ {\prime I}\tag{16}
\]

对一切 \(a' \in \mathbf{A}'\) 、\(b' \in \mathbf{B}'\) 、\(x \in \mathbf{E}\) 、\(y \in \mathbf{F}\) 。

\(\Phi'\) 的唯一性由张量积 \(a' \otimes x\) 与 \(b' \otimes y\) 分别生成 E′ 与 F′ 这一事实得出。为证其存在性，考虑映射

\[
m: (a, x, b ^ {\prime}, y) \rightarrow a ^ {\prime}. u (\Phi (x, y)). b ^ {\prime I}
\]

它把 \(\mathbf{A}' \times \mathbf{E} \times \mathbf{B}' \times \mathbf{F}\) 映到 \(\mathbf{G}\) 中。它显然是 \(\mathbf{Z}\)-重线性的，且把 (11) 与 (15) 考虑在内，满足 \(m(a', ax, b', y) = m(a'h(a), x, b', y)\) 与 \(m(a', x, b', by) = a'. u(\Phi(x, y)) \cdot h'(b)^{I} b'^{I} = m(a', x, b'h'(b), y)\) （ \(a \in \mathbf{A}, b \in \mathbf{B}, a' \in \mathbf{A}'\) ， \(b' \in \mathbf{B}'\) ， \(x \in \mathbf{E}\) ， \(y \in \mathbf{F}\) ）。因此存在一个 \(\mathbf{Z}\)-双线性映射 \(\Phi'\) （从 \(\mathbf{E}' \times \mathbf{F}'\) 到 \(\mathbf{G}'\) 中）

它满足 (16) （第 III 章，第 2 版，附录 II，第 1 小节，命题 2）。这一关系连同由 (12) 给出的 E′ 与 F′ 的模结构定义，并鉴于 (15)，表明 Φ′ 是关于 I 半双线性的，证毕。

\((\mathbf{A}^{\prime}, \mathbf{B}^{\prime})\) -双模 \(G^{\prime}\) （配备了满足 (11) 的、G 到 \(G^{\prime}\) 中的 Z-线性映射 u）的最重要例子如下：

\begin{enumerate}
\def\labelenumi{\arabic{enumi})}
\tightlist
\item
  我们取 G′ 为张量积 A′ ⊗\_A G ⊗\_B B′ （第 III 章，第 2 版，附录 II，第 9 小节），并取 u 为映射
\end{enumerate}

\[
g \to 1 \otimes g \otimes 1\tag{\(g\in G\)}
\]

它把 G 映到 G′ 中。如此定义的偶 (G′, u) 显然在下列意义下是泛的：对每个 (A′, B′)-双模 G₁′ 以及每个满足 (11) 的类似条件的、G 到 G₁′ 中的 Z-线性映射 u₁，存在唯一的 G′ 到 G₁′ 中的 Z-线性映射 f，使得 f(a'g'b') = a'f(g')b' （a′ ∈ A′, g′ ∈ G′, b′ ∈ B′；换言之，f 是 G′ 与 G₁′ 的双模结构之间的一个同态），且使得 u₁ = f ∘ u。

\begin{enumerate}
\def\labelenumi{\arabic{enumi})}
\setcounter{enumi}{1}
\item
  当 A = B = G（A 的 (A, A)-双模结构由左、右位似定义）、\(A' = B'\) 且 \(h = h'\) 时，可取 \(G'\) 为环 \(A'\) ，并取 u 为 A 到 \(A'\) 中的同态 h。
\item
  设我们有 A = B、\(A' = B'\) 、\(h = h'\) ，环 A 与 \(A'\) 都是交换的，且 G 的左 A-模结构与其右 A-模结构一致。这时可取 \(G'\) 为张量积 \(A' \otimes_{A} G\) （右 \(A'\) -模结构：即 \(G'\) 的这一结构与其左 \(A'\) -模结构相一致），并取 u 为映射 \(g \to 1 \otimes g\) （ \(g \in G\) ），它把 G 映到 \(G'\) 中。这时我们说，由命题 1（相应地命题 2）所定义的双线性（相应地半双线性）映射 \(\Phi'\) 是由 \(\Phi\) 经基环的扩张、亦即标量扩张而得到的。
\end{enumerate}

以下内容对双线性映射与半双线性映射同样有效；假设与记号即命题 1（相应地命题 2）者。给定 E 或 F 的一个子模 M，我们用 M′ 表示由 M 的典范像所生成的 E′ 或 F′ 的子模。

命题 3. --- 在命题 1（相应地命题 2）的假设与记号下，再设 A、B、A′、B′ 都是体，且映射 \(\alpha\) 与 \(\beta\) ——它们把 A′⊗A G 与 G⊗B B′ 映到 G′ 中，分别由 \(\alpha(a'\otimes g)=a'u(g)\) 与 \(\beta(g\otimes b')=u(g)b'\) （ \(a'\in A'\) ， \(b'\in B'\) ， \(g\in G\) ）刻画——都是单射。设 M 是 E 的一个子空间，N 是 F 的一个子空间。则子空间 \((M')^{0}\) （F′ 中关于 \(\Phi'\) 与 M′ 正交者）等于 \((M^{0})^{\prime}\) ，类似地有 \((N')^{0} = (N^{0})^{\prime}\) 。

事实上，包含关系 \((\mathbf{M}^{0})^{\prime}\subset(\mathbf{M}^{\prime})^{0}\) 与 \((\mathbf{N}^{0})^{\prime}\subset(\mathbf{N}^{\prime})^{0}\) 是显然的（而且对 A、B、A′、B′、\(\alpha\) 或 \(\beta\) 不作任何假设时也为真）。我们将证明包含关系 \((\mathbf{M}^{\prime})^{0}\subset(\mathbf{M}^{0})^{\prime}\) ；我们把完全类似的包含关系 \((\mathbf{N}^{\prime})^{0}\subset(\mathbf{N}^{0})^{\prime}\) 的验证留给读者。设 \(y^{\prime}\) 是 \((\mathbf{M}^{\prime})^{0}\) 的一个元素。可写

\[
y ^ {\prime} = \sum_ {i = 1} ^ {s} y _ {i} \otimes b _ {i} ^ {\prime} \quad (\text { resp. } y ^ {\prime} = \sum_ {i = 1} ^ {s} b _ {i} ^ {\prime} \otimes y _ {i})
\]

其中 \(y_i \in \mathbf{F} (1 \leqslant i \leqslant s)\) ，而诸 \(b_i'\) 是 \(\mathbf{B}'\) 的元素，在 \(\mathbf{B}\) 上线性无关（对于 \(\mathbf{B}\) 在 \(\mathbf{B}'\) 上的左（相应地右）模结构）。设 \(x \in \mathbf{M}\) 与 \(x' = 1 \otimes x \in \mathbf{M}'\) 。我们有

\[
0 = \Phi^ {\prime} (x ^ {\prime}, y ^ {\prime}) = \sum_ {i} u (\Phi (x, y _ {i})) b _ {i} ^ {\prime} = \beta (\sum_ {i} \Phi (x, y _ {i}) \otimes b _ {i} ^ {\prime})
\]

\[
(\text { resp. } 0 = \Phi^ {\prime} (x ^ {\prime}, y ^ {\prime}) = \sum_ {i} u (\Phi (x, y _ {i})) b _ {i} ^ {\prime \mathrm{I}} = \beta (\sum_ {i} \Phi (x, y _ {i}) \otimes b _ {i} ^ {\prime \mathrm{I}}).
\]

由于 β 是单射，且诸 \(b_i'\) （相应地诸 \(b _ {i} ^ {\prime \mathrm{I}}\) ，把 (15) 考虑在内）在 B′ 的左 B-模结构下于 B 上线性无关，这就导出 \(\Phi(x, y_i) = 0\) （对 \(i = 1, \ldots, s\)）。由于这最后一式对每个 \(x \in \mathbf{M}\) 都成立，故有 \(y_i \in \mathbf{M}^0\) （对 \(i = 1, \ldots, s\)），从而 \(y' \in (\mathbf{M}^0)'\) 。证毕。

推论. --- 在命题 3 的假设与记号下，\(\Phi'\) 非退化必须且只需 \(\Phi\) 非退化。

事实上，由命题 3，我们有 \((\mathbf{F}')^0 = (\mathbf{F}^0)'\) 与 \((\mathbf{E}')^0 = (\mathbf{E}^0)'\) 。另一方面，\(\Phi\) （相应地 \(\Phi'\) ）非退化必须且只需 \(\mathbf{F}^0 = \mathbf{E}^0 = \{0\}\) （相应地 \((\mathbf{F}')^0 = (\mathbf{E}')^0 = \{0\}\) ）。

注记. --- 设 A、B、A′ 与 B′ 都是体。那么，对上面三个例子中所定义的双模 G′，映射 α 与 β 都是单射，这由第 III 章，第 2 版，附录 II，第 6 小节立即得出。

